\subsection{Isotypical Components of a Module}

DEFINITION 4. --- Let A be a ring, M an A-module, and S a simple A-module. The isotypical component of type S of M, denoted by \(M_{S}\) , is the sum of the submodules of M isomorphic to S.

It is clear that \(M_{S}\) is the greatest submodule of M that is isotypical of type S. Because every submodule of \(M_{S}\) is isotypical of type S (VIII, p.~61, Proposition 2), we have \(N_{S} = M_{S} \cap N\) for every submodule N of M.

If \(S'\) is a simple A-module isomorphic to S, then we clearly have \(M_{S} = M_{S'}\) , so \(M_{S}\) depends only on the class of S (VIII, p.~51).

Let M be an A-module. There exists a greatest semisimple submodule of M, called the socle of M; it is the sum of the simple submodules of M and also the sum of the isotypical components of M. In particular, M is semisimple if and only if it is equal to its socle.

PROPOSITION 4. --- Let A be a ring. Denote the set of classes of simple A-modules by S. Let M be a semisimple A-module.

\begin{enumerate}
\def\labelenumi{\alph{enumi})}
\item
  The module M is the direct sum of the family \((\mathrm{M}_{\lambda})_{\lambda\in\mathcal{S}}\) of its isotypical components.
\item
  Suppose that M is the direct sum of a family \((\mathrm{N}_{i})_{i\in\mathrm{I}}\) of simple submodules. For every \(\lambda\in\mathscr{S}\) , let \(\mathrm{I}(\lambda)\) be the set of indices \(i\in I\) such that \(N_{i}\) is of class \(\lambda\) . We have \(M_{\lambda}=\bigoplus_{i\in\mathrm{I}(\lambda)}N_{i}\) .
\item
  If M is finitely generated, then the set of \(\lambda\in S\) such that \(M_{\lambda}\neq0\) is finite.
\item
  For every submodule \(\mathrm{N}\) of \(\mathrm{M}\) and every \(\lambda \in \mathcal{S}\) , we have \(\mathrm{N}_{\lambda} = \mathrm{N} \cap \mathrm{M}_{\lambda}\) and \((\mathrm{M} / \mathrm{N})_{\lambda} = (\mathrm{M}_{\lambda} + \mathrm{N}) / \mathrm{N}\) .
\end{enumerate}

Since M is semisimple, it is the sum of the family \((\mathrm{M}_{\lambda})_{\lambda\in\mathcal{S}}\) ; let us prove that this sum is direct. Let \(\lambda\in S\) . Denote the sum of the family \((\mathrm{M}_{\mu})_{\mu\in\mathcal{S}-\{\lambda\}}\) by \(M_{\lambda}^{\prime}\) . The module \(M_{\lambda}^{\prime}\) is the direct sum of a family of simple modules not isomorphic to \(\lambda\) (VIII, p.~56, Theorem 1). By Corollary 3 of Theorem 1 of VIII, p.~56, \(M_{\lambda}^{\prime}\) does not contain any simple submodules of class \(\lambda\) . We consequently have \(M_{\lambda}\cap M_{\lambda}^{\prime}=0\) . Assertion a) is therefore proved. By construction, we have \(M_{\lambda}\supset\bigoplus_{i\in\mathrm{I}(\lambda)}N_{i}\) , so assertion b) follows from Remark 1 of II, §1, No.~8, p.~208.

Assertion c) follows from a) and Proposition 23 of II, §1, No.~12, p.~221.

Let N be a submodule of M and \(\lambda\in S\) . The isotypical component \(N_{\lambda}\) of type \(\lambda\) of N is contained in \(M_{\lambda}\) and \(M_{\lambda}\cap N\subset N_{\lambda}\) . The intersection \(N\cap M_{\lambda}\) is therefore the isotypical component of N of type \(\lambda\) .

For every \(\lambda \in \mathcal{S}\) , the module \(\mathrm{M}_{\lambda} + \mathrm{N} / \mathrm{N}\) is isomorphic to \(\mathrm{M}_{\lambda} / (\mathrm{M}_{\lambda} \cap \mathrm{N})\) . It is therefore isotypical of type \(\lambda\) and contained in \((\mathrm{M} / \mathrm{N})_{\lambda}\) . The last assertion then follows from a) and II, §1, No.~8, p.~208, Remark 1.

COROLLARY. --- Let A be a ring and S the set of classes of simple A-modules. Let M be a semisimple A-module and N a submodule of M. Then we have \(N = \bigoplus_{\lambda \in \mathscr{S}} N \cap M_{\lambda}\) and \(M/N = \bigoplus_{\lambda \in \mathscr{S}} (M_{\lambda} + N)/N\) .

Since N and M/N are semisimple (VIII, p.~56, Corollary 3), the corollary follows from Proposition 4, d).

The support of a semisimple A-module M is the set of classes \(\lambda\) of simple A-modules such that the isotypical component of M of type \(\lambda\) is nonzero. The support of a finitely generated semisimple A-module is finite.

PROPOSITION 5. --- Let A be a ring, and let S be the set of classes of simple A-modules. Let M and N be A-modules.

\begin{enumerate}
\def\labelenumi{\alph{enumi})}
\item
  Let \(f: \mathrm{M} \to \mathrm{N}\) be a homomorphism. For every \(\lambda \in \mathcal{S}\) , \(f\) induces a homomorphism \(f_{\lambda}\) from \(\mathrm{M}_{\lambda}\) to \(\mathrm{N}_{\lambda}\) ; if \(\mathrm{M}\) is semisimple and \(f\) surjective, then each of the homomorphisms \(f_{\lambda}\) is surjective.
\item
  Suppose that M is semisimple. The mapping \(f \mapsto (f_{\lambda})_{\lambda \in \mathcal{S}}\) is a group isomorphism from \(\mathrm{Hom}_{\mathrm{A}}(\mathrm{M},\mathrm{N})\) to \(\prod_{\lambda \in \mathcal{S}}\mathrm{Hom}_{\mathrm{A}}(\mathrm{M}_{\lambda},\mathrm{N}_{\lambda})\) . When M is equal to N, the mapping is a ring isomorphism from \(\mathrm{End}_{\mathrm{A}}(\mathrm{M})\) to \(\prod_{\lambda \in \mathcal{S}}\mathrm{End}_{\mathrm{A}}(\mathrm{M}_{\lambda})\) .
\end{enumerate}

For every \(\lambda \in \mathcal{S}\) , the submodule \(f(\mathrm{M}_{\lambda})\) of \(\mathbf{N}\) is isomorphic to a quotient of an isotypical module of type \(\lambda\) ; it is therefore isotypical of type \(\lambda\) and, consequently, contained in \(\mathrm{N}_{\lambda}\) .

Suppose that \(\mathrm{M}\) is semisimple and \(f\) surjective. Then \(f\) induces an isomorphism from \(\mathrm{M} / \mathrm{Ker}(f)\) to \(\mathrm{N}\) that sends \((\mathrm{M}_{\lambda} + \mathrm{Ker}(f)) / \mathrm{Ker}(f)\) to \(f(\mathrm{M}_{\lambda})\) .

By Proposition 4 of VIII, p.~65, we have \(N_{\lambda} = f(M_{\lambda})\) , which completes the proof of a).

The mapping considered in b) is clearly a group homomorphism, and it is a ring homomorphism when M is equal to N. Let \((f_{\lambda})_{\lambda \in \mathcal{S}}\) be an element of \(\prod_{\lambda \in \mathcal{S}} \mathrm{Hom}(M_{\lambda}, N_{\lambda})\) . The unique element of its inverse image under the mapping in b) is the homomorphism \(f : M \to N\) defined by

\[
f \Bigl (\sum_ {\lambda \in \mathcal {S}} x _ {\lambda} \Bigr) = \sum_ {\lambda \in \mathcal {S}} f _ {\lambda} (x _ {\lambda})
\]

for every \((x_{\lambda})_{\lambda \in \mathcal{S}}\in \bigoplus_{\lambda}\mathrm{M}_{\lambda}\).

Remark. --- Let A and B be rings. Let M be an (A, B)-bimodule. It follows from Proposition 5 that the isotypical components of the A-module M are sub-bimodules of M. This holds, in particular, when M is an A-module and B is the opposite ring of \(\operatorname{End}_{\mathrm{A}}(\mathrm{M})\) .

Example. --- Let us consider the case when the ring A is commutative. The mapping that sends a maximal ideal m onto \(\operatorname{cl}(A/m)\) is a bijection from the set of maximal ideals of A to the set S of classes of simple A-modules (VIII, p.~51). The inverse bijection sends \(\lambda\) to its annihilator \(m_{\lambda}\) .

Let M be an A-module. For every \(\lambda\in S\) , the isotypical component \(M_{\lambda}\) of type \(\lambda\) of M consists of the elements annihilated by \(m_{\lambda}\) , and we can view \(M_{\lambda}\) as a vector space over the field \(A/m_{\lambda}\) . If M is semisimple and N is another A-module, then we can deduce a group isomorphism from \(\operatorname{Hom}_{\mathrm{A}}(\mathrm{M},\mathrm{N})\) to \(\prod_{\lambda\in\mathcal{S}}\operatorname{Hom}_{\mathrm{A}/m_{\lambda}}(\mathrm{M}_{\lambda},\mathrm{N}_{\lambda})\) from Proposition 5.

COROLLARY 1. --- Let M be a semisimple A-module and N a submodule of M. The following properties are equivalent:

\begin{enumerate}
\def\labelenumi{(\roman{enumi})}
\item
  There exists a unique submodule supplementary to N in M.
\item
  We have \(\mathrm{Hom}_{\mathrm{A}}(\mathrm{M} / \mathrm{N},\mathrm{N}) = 0\)
\item
  There exists a subset \(\Lambda\) of \(\mathcal{S}\) such that \(N = \bigoplus_{\lambda \in \Lambda} M_{\lambda}\) .
\end{enumerate}

Choose a submodule \(N^{\prime}\) supplementary to N in M (VIII, p.~56, Corollary 2). If we identify M with \(N^{\prime} \times N\) , then the submodules of M supplementary to N are the graphs of the A-linear mappings from \(N^{\prime}\) to N. Since \(N^{\prime}\) is isomorphic to M/N, we have proved the equivalence of properties (i) and (ii). By Proposition 5, b), the group \(\operatorname{Hom}_{\mathrm{A}}(\mathrm{N}^{\prime}, \mathrm{N})\) is isomorphic to the group \(\prod_{\lambda \in \mathcal{S}} \operatorname{Hom}_{\mathrm{A}}(\mathrm{N}_{\lambda}^{\prime}, \mathrm{N}_{\lambda})\) . It is zero if and only if for every \(\lambda \in S\) , we have \(N_{\lambda} = 0\) or \(N_{\lambda}^{\prime} = 0\) (VIII, p.~61, Remark), that is, \(N_{\lambda} = 0\) or \(N_{\lambda} = M_{\lambda}\) . This proves the equivalence of properties (ii) and (iii).

COROLLARY 2. --- Let M be an A-module. The following two conditions are equivalent:

\begin{enumerate}
\def\labelenumi{(\roman{enumi})}
\item
  Every submodule of M admits a unique supplementary submodule.
\item
  The module M is the direct sum of a family \((\mathrm{S}_{i})_{i\in\mathrm{I}}\) of simple, pairwise nonisomorphic modules.
\end{enumerate}

Suppose that M satisfies these conditions. Then, for every submodule N of M, there exists a unique subset J of I such that we have \(N = \bigoplus_{j \in J} S_j\) and every simple submodule of M is equal to one of the \(S_i\) .

Conditions (i) and (ii) both imply that M is semisimple.

Suppose that condition (i) is satisfied. Let \(\lambda \in \mathcal{S}\) . By the equivalence of (i) and (iii) in Corollary 1, every submodule of \(\mathrm{M}_{\lambda}\) is zero or equal to \(\mathrm{M}_{\lambda}\) . Consequently, \(\mathrm{M}_{\lambda}\) is zero or simple, and (ii) follows from the fact that \(\mathrm{M}\) is the direct sum of the family \((\mathrm{M}_{\lambda})_{\lambda \in \mathcal{S}}\) .

Conversely, if condition (ii) is satisfied, then \(M_{\lambda}\) is zero or simple for every \(\lambda \in S\) . If N is a submodule of M, then \(N \cap M_{\lambda}\) is equal to 0 or \(M_{\lambda}\) for every \(\lambda \in S\) . Since we have \(N = \bigoplus_{\lambda \in S} (N \cap M_{\lambda})\) (VIII, p.~66, Corollary), the submodule N has property (iii) of Corollary 1 and therefore admits a unique supplementary submodule in M. This proves that (ii) implies (i), as well as the last assertions of the corollary.

Let M be an A-module and S a simple A-module. Denote the opposite ring of the field \(\operatorname{End}_{\mathrm{A}}(\mathrm{S})\) by D, and view S as an (A,D)-bimodule. Then \(\operatorname{Hom}_{\mathrm{A}}(\mathrm{S},\mathrm{M})\) is a left vector space over D, and \(\operatorname{Hom}_{\mathrm{A}}(\mathrm{M},\mathrm{S})\) is a right vector space over D. The dual of the left D-vector space \(\operatorname{Hom}_{\mathrm{A}}(\mathrm{S},\mathrm{M})\) is a right vector space over D (II, §2, No.~3, p.~232, Definition 2). For every \(u \in \operatorname{Hom}_{\mathrm{A}}(\mathrm{M},\mathrm{S})\) , the mapping \(h(u): v \mapsto u \circ v\) from \(\operatorname{Hom}_{\mathrm{A}}(\mathrm{S},\mathrm{M})\) to \(\operatorname{Hom}_{\mathrm{A}}(\mathrm{S},\mathrm{S}) = \mathrm{D}\) is a linear form on the left D-vector space \(\operatorname{Hom}_{\mathrm{A}}(\mathrm{S},\mathrm{M})\) .

PROPOSITION 6. --- Keep the notation above, and suppose that the A-module M is semisimple. The mapping \(u \mapsto h(u)\) from the right D-vector space \(\mathrm{Hom}_{\mathrm{A}}(\mathrm{M},\mathrm{S})\) to the dual of the left D-vector space \(\mathrm{Hom}_{\mathrm{A}}(\mathrm{S},\mathrm{M})\) is D-linear and bijective.

Let \(u \in \mathrm{Hom}_{\mathrm{A}}(\mathrm{M}, \mathrm{S})\) , \(v \in \mathrm{Hom}_{\mathrm{A}}(\mathrm{S}, \mathrm{M})\) , and \(d \in \mathrm{D}\) . We have

\[
h (u d) (v) = h (d \circ u) (v) = d \circ u \circ v = d \circ (h (u) (v)) = h (u) (v) d.
\]

This proves that the mapping h is D-linear. It is simply the mapping given by \(u \mapsto \operatorname{Hom}(1_{\mathrm{S}}, u)\) from \(\operatorname{Hom}_{\mathrm{A}}(\mathrm{M}, \mathrm{S})\) to \(\operatorname{Hom}_{\mathrm{D}}(\operatorname{Hom}_{\mathrm{A}}(\mathrm{S}, \mathrm{M}), \operatorname{Hom}_{\mathrm{A}}(\mathrm{S}, \mathrm{S}))\) . To prove that it is bijective, by Proposition 5, b) of VIII, p.~66, it suffices to treat the case when M is isotypical of type S; we can then apply the corollary of VIII, p.~65.

